%\documentclass[landscape,a0b,final,a4resizeable]{a0poster}
\documentclass[landscape,a0b,final]{a0poster}
%\documentclass[portrait,a0b,final,a4resizeable]{a0poster}
%\documentclass[portrait,a0b,final]{a0poster}
%%% Option "a4resizeable" makes it possible ot resize the
%   poster by the command: psresize -pa4 poster.ps poster-a4.ps
%   For final printing, please remove option "a4resizeable" !!

\usepackage{epsfig}
\usepackage{multicol}
\usepackage{pstricks,pst-grad}
\usepackage{braket}

%%%%%%%%%%%%%%%%%%%%%%%%%%%%%%%%%%%%%%%%%%%
% Definition of some variables and colors
%\renewcommand{\rho}{\varrho}
%\renewcommand{\phi}{\varphi}
\setlength{\columnsep}{3cm}
\setlength{\columnseprule}{2mm}
\setlength{\parindent}{0.0cm}

%My commands
\usepackage{overcite}
\usepackage{amsmath}
\usepackage{amsfonts}
\usepackage{bm}
\newcommand{\mat}[1]{\mathbb{#1}} % mtpro special fonts
\newcommand{\D}{\ensuremath{\partial }} %% partial derivative
\newcommand{\tensortwo}[1]{\mat{#1}}
\newcommand{\bindent}{\ \ \ \ }
\newcommand{\vect}[1]{\bm{#1}} %% boldface vectors
\newcommand{\feb}[1]{\ensuremath{\beta^i_{#1 #1} (\vect{r})}}
\renewcommand{\Im}{\operatorname{Im}} %% Re and Im as it is the default

%%%%%%%%%%%%%%%%%%%%%%%%%%%%%%%%%%%%%%%%%%%%%%%%%%%%
%%%               Background                     %%%
%%%%%%%%%%%%%%%%%%%%%%%%%%%%%%%%%%%%%%%%%%%%%%%%%%%%

\newcommand{\background}[3]{
  \newrgbcolor{cgradbegin}{#1}
  \newrgbcolor{cgradend}{#2}
  \psframe[fillstyle=gradient,gradend=cgradend,
  gradbegin=cgradbegin,gradmidpoint=#3](0.,0.)(1.\textwidth,-1.\textheight)
}

%%%%%%%%%%%%%%%%%%%%%%%%%%%%%%%%%%%%%%%%%%%%%%%%%%%%
%%%                Poster                        %%%
%%%%%%%%%%%%%%%%%%%%%%%%%%%%%%%%%%%%%%%%%%%%%%%%%%%%

\newenvironment{poster}{
  \begin{center}
  \begin{minipage}[c]{0.981\textwidth}
}{
  \end{minipage} 
  \end{center}
}

%%%%%%%%%%%%%%%%%%%%%%%%%%%%%%%%%%%%%%%%%%%%%%%%%%%%
%%%                pcolumn                       %%%
%%%%%%%%%%%%%%%%%%%%%%%%%%%%%%%%%%%%%%%%%%%%%%%%%%%%

\newenvironment{pcolumn}[1]{
  \begin{minipage}{#1\textwidth}
  \begin{center}
}{
  \end{center}
  \end{minipage}
}

%%%%%%%%%%%%%%%%%%%%%%%%%%%%%%%%%%%%%%%%%%%%%%%%%%%%
%%%                pbox                          %%%
%%%%%%%%%%%%%%%%%%%%%%%%%%%%%%%%%%%%%%%%%%%%%%%%%%%%

\newrgbcolor{lcolor}{0. 0. 0.80}
\newrgbcolor{gcolor1}{1. 1. 1.}
\newrgbcolor{gcolor2}{.80 .80 1.}

\newcommand{\pbox}[4]{
\psshadowbox[#3]{
\begin{minipage}[t][#2][t]{#1}
#4
\end{minipage}
}}

%%%%%%%%%%%%%%%%%%%%%%%%%%%%%%%%%%%%%%%%%%%%%%%%%%%%
%%%                myfig                         %%%
%%%%%%%%%%%%%%%%%%%%%%%%%%%%%%%%%%%%%%%%%%%%%%%%%%%%
% \myfig - replacement for \figure
% necessary, since in multicol-environment 
% \figure won't work

\newcommand{\myfig}[3][0]{
\begin{center}
  \vspace{0.25cm}
  \includegraphics[width=#3\hsize,angle=#1]{#2}
  \nobreak\medskip
\end{center}}

%%%%%%%%%%%%%%%%%%%%%%%%%%%%%%%%%%%%%%%%%%%%%%%%%%%%
%%%                mycaption                     %%%
%%%%%%%%%%%%%%%%%%%%%%%%%%%%%%%%%%%%%%%%%%%%%%%%%%%%
% \mycaption - replacement for \caption
% necessary, since in multicol-environment \figure and
% therefore \caption won't work

%\newcounter{figure}
\setcounter{figure}{1}
\newcommand{\mycaption}[1]{
  \vspace{0.25cm}
  \begin{quote}
    {{\sc Figure} \arabic{figure}: #1}
  \end{quote}
  \vspace{0.25cm}
  \stepcounter{figure}
}

%%%%%%%%%%%%%%%%%%%%%%%%%%%%%%%%%%%%%%%%%%%%%%%%%%%%%%%%%%%%%%%%%%%%%%
%%% Begin of Document
%%%%%%%%%%%%%%%%%%%%%%%%%%%%%%%%%%%%%%%%%%%%%%%%%%%%%%%%%%%%%%%%%%%%%%

\begin{document}

\newrgbcolor{lightblue}{0. 0. 0.80}
\newrgbcolor{white}{1. 1. 1.}
\newrgbcolor{whiteblue}{.80 .80 1.}
\newrgbcolor{ubblue}{0.0157 0.102 0.608}
\newrgbcolor{ubgray}{0.4 0.4 0.4}

\background{1. 1. 1.}{1. 1. 1.}{0.5}

\vspace*{0.25cm} %Controls white space on top of poster

\begin{poster}

%%%%%%%%%%%%%%%%%%%%%
%%% Header
%%%%%%%%%%%%%%%%%%%%%
\begin{center}
\begin{pcolumn}{0.98}

\pbox{0.98\textwidth}{}{linewidth=2mm,framearc=0.3,linecolor=ubblue,fillstyle=gradient,gradangle=0,gradbegin=ubgray,gradend=white,gradmidpoint=1.0,framesep=1em}{
    %%% Unisiegel
    \begin{minipage}[c][9cm][c]{0.1\linewidth}
    \begin{center}
        \includegraphics[width=20cm,angle=0]{figures/ublogo}
    \end{center}
    \end{minipage}
    %%% Titel
    \begin{minipage}[c][9cm][c]{0.75\textwidth}
    \begin{center}
        %hspace commands control centering of title above poster. You may need to mess with these depending on how long your title is.
        \hspace{5cm}
        {\sc \huge Analysis of Optical Activity in Terms of Bonds and Lone-Pairs:}\\[5mm] 
        \hspace{5cm}
        {\sc \huge The Exceptionally Large Optical Rotation of Norborneneone\cite{Autschbach:2012k}}\\[5mm]
        \hspace{5cm}
        {\Large Barry Moore II, Monika Srebro, and Jochen Autschbach\\[7.5mm]
        \hspace{5cm}
        University at Buffalo, State University of New York, Buffalo, NY, 14260}
    \end{center}
    \end{minipage}
    %%% GK-Logo
    \begin{minipage}[c][9cm][c]{0.01\linewidth}
    \begin{flushright}
        %This hspace puts the nsf symbol at the far right of the poster banner. May need to change this but don't remove the '*'
        \hspace*{5cm}
        \includegraphics[width=10cm,angle=0]{figures/nsf4}
    \end{flushright}
    \end{minipage}
}
\end{pcolumn}
\end{center}

\vspace*{0.05cm} %Spacing between banner and text

%%%%%%%%%%%%%%%%%%%%%
%%% Content
%%%%%%%%%%%%%%%%%%%%%

%%% Begin of Multicols-Enviroment
\begin{multicols}{3}

%%% Introduction
\begin{center}\pbox{0.8\columnwidth}{}{linewidth=2mm,framearc=0.3,linecolor=ubblue,fillstyle=gradient,gradangle=0,gradbegin=white,gradend=ubgray,gradmidpoint=1.0,framesep=1em}{\begin{center} \Large \color{ubblue}{\bf{Introduction}}\end{center}}\end{center} %Makes the Abstract banner 
\vspace{0.25cm}
    \large    
    \begin{center}
    \myfig[0]{figures/structures-sos}{1.0}
    \mycaption{a) Some bicyclic organic molecules, and measured specific
    and molar optical rotations (in usual units) for $\lambda = 589.3$
    nm. b) SOS molar optical rotation of norbornenone
    (PBE0/aug-cc-pVDZ).  Experimental solution-phase OR and the result
    from a direct calculation of the OR.}
    \end{center}

%%% Functional Tuning
\vspace{0.25cm}
\begin{center}\pbox{0.8\columnwidth}{}{linewidth=2mm,framearc=0.3,linecolor=ubblue,fillstyle=gradient,gradangle=0,gradbegin=white,gradend=ubgray,gradmidpoint=1.0,framesep=1em}{\begin{center} \Large \color{ubblue}{\bf{Functional Tuning}}\end{center}}\end{center}
\vspace{0.25cm}

    \bindent We use the following 3-parameter expression for the inverse
    interelectronic distance in the exchange functional
    \cite{yanai_new_2004}:
    \begin{equation}
    \label{eq:r12-range-separation-alpha-beta-gamma}
    \frac{1}{r_{12}} =
    \frac{\alpha + \beta \operatorname{erf} (\gamma r_{12})}{r_{12}} + 
    \frac{1-[\alpha + \beta \operatorname{erf} (\gamma r_{12})]}{r_{12}} 
    \end{equation}
    where the first term on the right-hand side is used for the long-range
    HF part. In Equation (\ref{eq:r12-range-separation-alpha-beta-gamma}),
    $\gamma$ is the range-separation parameter, $\alpha$ is the HF
    fraction at very small interelectronic distances, and $\alpha+\beta$
    quantifies the fraction of HF exchange in the asymptotic limit at
    large interelectronic separations. A global hybrid such as B3LYP
    corresponds to $\beta = 0$, with $\alpha$ being the HF fraction of
    exchange in the functional.

    \bindent We consider a range-separated hybrid variant of the PBE functional
    which satisfies the correct asymptotic behavior of the XC
    potential. It follows that $\alpha + \beta = 1$. Further, the
    functional is tuned to satisfy
    $-{\varepsilon}^{\text{HOMO}}(N)=\text{IP}(N)$ subject to the $\alpha
    + \beta = 1$ constraint. Following Reference
    \citen{kuritz_charge-transfer-like_2011}, we minimize
    \begin{equation}
    J^2 =
    \sum_{i=0}^{1} [\varepsilon^{\text{HOMO}}(N+i)+\text{IP}(N+i)]^2 
    \label{eq:J2-criterion}
    \end{equation} 
    to satisfy the IP condition simultaneously for the system with $N$ and
    $N+1$ electrons, where $\text{IP}(N)= E(N)-E(N-1)$ is calculated as a
    difference of total energies. This optimization is performed in the
    $\alpha = 1-\beta$ vs.\ $\gamma$ plane, which is the reason why the
    procedure is referred to as a `two-dimensional' tuning here. The
    tuning was accomplished as follows. Hartree-Fock molecular orbitals
    were calculated for the neutral ($N$ electrons), anion ($N+1$), and
    cation ($N-1$) systems. These molecular orbital vectors were then used
    as initial guesses for DFT calculations with the PBE0 global hybrid
    functional. The PBE0 molecular orbital vectors were used as input for
    calculations with $\alpha$ between 0 and 0.5 and $\gamma$ between 0
    and 0.5 on the tuning grid. MOs from calculations with $\alpha = 0.5$
    and a given $\gamma$ were used as initial guesses for calculations
    with $\alpha$ ranging from 0.6 to 0.8 and the same value of $\gamma$
    in order to improve convergence of the self-consistent field (SCF)
    procedure.  Finally, for sets $\{\alpha,\beta = 1-\alpha,\gamma\}$
    satisfying the IP condition, we investigate the DFT delocalization
    error, quantified by the curvature of $E(N)$.
    %%% Figures:
    \begin{center}
    \myfig[0]{figures/norbornenone-tuning-correct}{1.0}
    \mycaption{ a) Contour plot of $J^2$ as a function of
    $\alpha$ and $\gamma$. The thick orange line corresponds to $J^2
    \approx 0$. b) Energy of norbornenone as a function
    of fractional electron number, ${\Delta}N$, relative to neutral
    system $({\Delta}N = 0$) for tuned LC-PBE0-\# parameterizations. The
    numerical values in parentheses correspond to $({\Delta}N)^2$
    coefficients of quadratic fits to $E(N)$ in the electron-deficient
    and electron-rich regime, respectively $({\Delta}N < 0$, ${\Delta}N
    > 0)$. The inset shows the $E(N)$ behavior in the electron-rich
    regime.}
    \end{center}

%%% Analysis Protocol
\vspace{1.00cm} %Increased spacer size normally 0.25cm
\begin{center}
\pbox{0.8\columnwidth}{}{linewidth=2mm,framearc=0.1,linecolor=ubblue,fillstyle=gradient,gradangle=0,gradbegin=white,gradend=ubgray,gradmidpoint=1.0,framesep=1em}
{\begin{center} \Large \color{ubblue}{\bf{Analysis Protocol}}\end{center}}
\end{center}
\vspace{1.00cm} %Increased spacer size normally 0.25cm
    %%% Figures:
    \begin{center}
        % \myfig[rotation in degrees]{epsfile}{stretching parameter}
    \myfig[0]{figures/flowchart}{1.0}
    \mycaption{Flowchart for OR analysis. CNC = center of nuclear charges,
    PAS = principal axis system of $\tilde{\tensortwo{B}}$.}
    \end{center}

    \begin{equation}
    \normalsize
    \label{eq:amos-eq-tddft}
    \beta(\omega) = - \frac{2}{3} \Im \sum_{u = x,y,z} \sum_i
    \Braket{\frac{\D \varphi_i}{\D E_u(+\omega)} + \frac{\D \varphi_i}{\D
    E_u(-\omega)}|\frac{\D \varphi_i}{\D B_u (0)}} 
    := \frac{1}{3} \sum_{u = x,y,z} \sum_i \int \feb{u} dV
    \end{equation}
\vspace{1cm} %SPECIAL SPACER
%%% Results and Discussion
\vspace{0.25cm}
\begin{center}
\pbox{0.8\columnwidth}{}{linewidth=2mm,framearc=0.1,linecolor=ubblue,fillstyle=gradient,gradangle=0,gradbegin=white,gradend=ubgray,gradmidpoint=1.0,framesep=1em}
{\begin{center} \Large \color{ubblue}{\bf{Results and Discussion}}\end{center}}
\end{center}
\vspace{0.25cm}
    
    %%% Figures:
    \begin{center}
        % \myfig[rotation in degrees]{epsfile}{stretching parameter}
    \myfig[0]{figures/scheme-pi-lp-comparison}{1.0}
    \mycaption{Left: Qualitative interpretation of the electric- and
    magnetic-field perturbations of the $\pi$LP(O) orbital at the
    carbonyl group via construction of the integrand of
    Eq.\ (\ref{eq:amos-eq-tddft}). Right: Static \feb{3}\ for $i
    =\pi$LP(O) (isosurfaces 0.028 au). The numbers listed are
    the total orbital contributions to $\beta_{33}$ obtained upon
    volume integration (italic), and approximate numerical
    integrations within low-value isosurfaces centered at the C=C and
    C=O moieties, respectively.}
    \end{center}
    
    %%% Figures:
    \begin{center}
        % \myfig[rotation in degrees]{epsfile}{stretching parameter}
    \myfig[0]{figures/norbornenone-pinene-comparison}{1.0}
    \mycaption{Top: Isosurfaces of the $\pi$(C=C) LMO (0.1 au),
    and their contributions to $\tensortwo{\beta}$ tensor for
    norbornenone and $\alpha$-pinene. Bottom: Isosurfaces (0.04 au) of
    \feb{1} and \feb{3} for $\pi$(C=C), and approximate numerical
    integrations within isosurfaces centered at the C=C and C=O
    moieties, respectively. LC-PBE0* functional.}
    \end{center}

%%% Acknowledgements
\vspace{0.25cm}
\begin{center}
\pbox{0.8\columnwidth}{}{linewidth=2mm,framearc=0.1,linecolor=ubblue,fillstyle=gradient,gradangle=0,gradbegin=white,gradend=ubgray,gradmidpoint=1.0,framesep=1em}
{\begin{center} \Large \color{ubblue}{\bf{Acknowledgements}}\end{center}}
\end{center}
\vspace{0.25cm}
    
     This research has been supported by the National Science
     Foundation, Grant No. CHE 0952253 to J.A. M.S. is grateful
     for financial support from the Foundation for Polish Science
     (START scholarship) as well as from the Polish Ministry of
     Science and Higher Education (Mobility Plus program). We
     thank the Center for Computational Research at the University
     at Buffalo for supporting our research.
      
%%% References
\bibliographystyle{achemso}
\bibliography{poster}

\end{multicols}
\end{poster}
\end{document}

